\documentclass[11pt]{article}
\usepackage[T1]{fontenc}
\usepackage{lmodern}
\usepackage[margin=1in]{geometry}
\usepackage{amsmath,amssymb,amsthm,mathtools}
\usepackage{enumitem}
\usepackage{microtype}
\usepackage[colorlinks=true,linkcolor=blue,citecolor=blue,urlcolor=blue]{hyperref}

\newcommand{\BL}{\mathcal{BL}}
\newcommand{\C}{\mathcal C}
\newcommand{\U}{\mathcal U}
\newcommand{\R}{\mathbb R}
\newcommand{\1}{\mathbf 1}
\newcommand{\supp}{\operatorname{supp}}
\newcommand{\dist}{\operatorname{dist}}
\newcommand{\dens}{\operatorname{dens}}
\newcommand{\Ran}{\operatorname{Ran}}
\newcommand{\Id}{\operatorname{Id}}

\newtheorem{theorem}{Theorem}[section]
\newtheorem{proposition}[theorem]{Proposition}
\newtheorem{lemma}[theorem]{Lemma}
\newtheorem{corollary}[theorem]{Corollary}
\newtheorem{claim}[theorem]{Claim}
\theoremstyle{definition}
\newtheorem{definition}[theorem]{Definition}
\theoremstyle{remark}
\newtheorem{remark}[theorem]{Remark}

\title{The real Hilbert endpoint in Raynaud's\\
$\mathcal{BL}_1L_q$ axiomatizability problem}
\author{Proof draft for independent verification}
\date{June 18, 2026}

\begin{document}
\maketitle

\begin{abstract}
Raynaud asked whether the Banach-space reduct of the mixed-norm Banach-lattice
class $\BL_1L_q$ is axiomatizable for $1<q\leq2$.  This note proves the
remaining real Hilbert endpoint $q=2$.

The main new device is a canonical radial truncation operation $\beta$ on
$L_1$-direct integrals of real Hilbert spaces.  Although the lattice
structure is not determined by the Banach norm at $q=2$, the operation
$\beta$ is: it admits an intrinsic description using the Boolean algebra of
$L$-projections and is therefore preserved by every surjective linear
isometry.  It also commutes with ultraproducts.  A Keisler--Shelah common-ultrapower/equalizer argument then reduces
ultraroot closure to a geometric statement about closed
$\beta$-invariant subspaces.  Such a subspace is shown to be an
$\ell_1$-sum of Hilbert-valued $L_1$-spaces by extracting a scalar
sub-$\sigma$-algebra and applying Raynaud's module representation lemma.
Consequently the real Banach-space class underlying $\BL_1L_2$ is closed
under ultraproducts and ultraroots, hence is axiomatizable.

A targeted literature search through June 18, 2026 did not locate this
endpoint argument in print.  The note is therefore presented as a complete
proof draft for independent expert verification before citation as an
established result.
\end{abstract}

\section{Statement of the result}

For $1\leq p,q<\infty$, let $\BL_pL_q$ denote the class of Banach lattices
lattice-isometric to bands in mixed-norm lattices $L_p(\Omega;L_q(S))$.
If $\mathcal A$ is a class of Banach lattices, write $\mathcal A^B$ for the
class of their underlying Banach spaces, up to surjective linear isometry.
Raynaud proved many Banach-space axiomatizability results for these classes
and left the following range open \cite[Question~6.17]{Raynaud2018}:
\[
       (\BL_1L_q)^B,\qquad 1<q\leq2.
\]
The purpose of this note is to settle the real endpoint.

\begin{theorem}[Main theorem]
\label{thm:main}
The class of real Banach spaces linearly isometric to members of
$\BL_1L_2$ is axiomatizable in the language of real Banach spaces.
Equivalently,
\[
             \boxed{(\BL_1L_2)^B\text{ is closed under ultraroots}.}
\]
\end{theorem}

Closure under ultraproducts was already known from the Banach-lattice theory
of $\BL_pL_q$ \cite{HensonRaynaud2007,Raynaud2018}.  Thus only ultraroot
closure needs proof.

The argument does not reconstruct the coordinate lattice in a Hilbert
fibre.  That would be impossible: arbitrary orthogonal operators preserve
the norm and need not preserve lattice disjointness.  Instead, it recovers a
weaker radial operation which is exactly strong enough to recognize the
underlying Banach-space class.

\section{The endpoint class and random norms}

Let $H$ be a real Hilbert space and let $(\Omega,\Sigma,\mu)$ be a localizable
measure space.  For $x\in L_1(\Omega;H)$ write
\[
             N(x)(\omega)=\|x(\omega)\|_H
\]
for its random norm.  We use the same notation in bands and direct
integrals.

Define $\C_1^{\R}$ to be the class of real Banach spaces isometric to spaces
of the form
\begin{equation}
\label{eq:C1}
       \left(\bigoplus_{\gamma\in\Gamma}
       L_1(\Omega_\gamma,\mu_\gamma;H_\gamma)\right)_1,
\end{equation}
where the $H_\gamma$ are real Hilbert spaces and the bases are localizable.

\begin{proposition}
\label{prop:C1=BL}
One has
\[
                   \C_1^{\R}=(\BL_1L_2)^B.
\]
Moreover every member of $\C_1^{\R}$ can be represented as a band, hence as
a closed $L_\infty(\Sigma)$-submodule, of some $L_1(\Omega;H)$ with a fixed
real Hilbert space $H$.
\end{proposition}

\begin{proof}
Using the standard localizable realization of the base, a band in
$L_1(\Omega;L_2(S))$ is a closed $L_\infty(\Sigma)$-submodule of
$L_1(\Omega;H)$, where $H=L_2(S)$, and its random norms are
$\Sigma$-measurable.  Raynaud's module representation lemma
\cite[Lemma~4.4]{Raynaud2004} (with $p=1$ and $\mathcal F=\Sigma$) gives a
random-norm-preserving isometry onto an $\ell_1$-sum of spaces
$L_1(\Omega_\gamma;H_\gamma)$.

Conversely, choose an orthonormal basis in every $H_\gamma$, identify it
with a coordinate band of one sufficiently large Hilbert lattice
$\ell_2(I)$, and take the disjoint union of the base spaces.  The space in
\eqref{eq:C1} is then a band in $L_1(\Omega;\ell_2(I))$.
\end{proof}

We shall also use the mixed-norm ultraproduct theorem of L\'evy--Raynaud and
Henson--Raynaud \cite{LevyRaynaud1984,HensonRaynaud2007,HensonRaynaud2011}.
In the form needed here it says the following.

\begin{proposition}[Random norms in ultraproducts]
\label{prop:random-ultra}
Let $(Z_i)_{i\in I}$ be members of $\C_1^{\R}$, each equipped with one of its
$\BL_1L_2$ lattice representations and random norm $N_i$.  For an
ultrafilter $\U$ on $I$, the Banach ultraproduct
$Z=\prod_{i,\U}Z_i$ admits a $\BL_1L_2$ representation with random norm $N$
satisfying
\[
             N([x_i]_{\U})=[N_i(x_i)]_{\U}.
\]
In particular, $\C_1^{\R}$ is closed under Banach ultraproducts.
\end{proposition}

\section{A canonical radial operation}

For a direct-integral representation $Z=L_1(\Omega;H_\omega)$, define
\begin{equation}
\label{eq:beta-def}
 \beta_Z(x,y)(\omega)=
 \begin{cases}
 \displaystyle
 \frac{\min\{N(x)(\omega),N(y)(\omega)\}}{N(y)(\omega)}\,y(\omega),
      &N(y)(\omega)>0,\\[1.2ex]
 0,   &N(y)(\omega)=0.
 \end{cases}
\end{equation}
Thus $\beta_Z(x,y)(\omega)$ is the radial projection of $y(\omega)$ onto the
Hilbert ball of radius $N(x)(\omega)$.  We next prove that this operation is
intrinsic to the Banach norm.

Recall that an \emph{$L$-projection} on a Banach space $Z$ is a bounded
projection $P$ such that
\[
          \|z\|=\|Pz\|+\|(I-P)z\|\qquad(z\in Z).
\]

\begin{lemma}[The base is intrinsic]
\label{lem:L-projections}
Let $Z=L_1(\Omega;H_\omega)$ be an $L_1$-direct integral of nonzero real
Hilbert spaces.  Its $L$-projections are exactly the base multipliers
\[
                     P_Az=\1_Az\qquad(A\in\Sigma),
\]
modulo null sets.
\end{lemma}

\begin{proof}
Every $P_A$ is plainly an $L$-projection.  Conversely, let $P$ be an
$L$-projection, $R=\Ran P$, and $K=\ker P$.  For $r\in R$ and $k\in K$,
\[
 \|r+k\|=\|r\|+\|k\|,
 \qquad
 \|r-k\|=\|r\|+\|k\|.
\]
Both pointwise triangle defects are nonnegative and have integral zero.
Hence both triangle equalities hold almost everywhere in every Hilbert
fibre.  Simultaneous equality for $r(\omega)+k(\omega)$ and
$r(\omega)-k(\omega)$ forces $r(\omega)=0$ or $k(\omega)=0$.

In the complete measure algebra, let $A$ be the supremum of the supports of
all elements of $R$.  Every element of $R$ is supported in $A$, and the
preceding base-disjointness forces every element of $K$ to be supported in
$A^c$.  If $z=Pz+(I-P)z$, it follows that $Pz=\1_Az$.
\end{proof}

For $L$-projections $Q,P$, write $Q\leq P$ when $\Ran Q\subseteq\Ran P$.
For $x,y\in Z$ and an $L$-projection $P$, put
\begin{equation}
\label{eq:meet-functional}
 m_{x,y}(P)=
 \inf_{Q\leq P}\bigl(\|Qx\|+\|(P-Q)y\|\bigr).
\end{equation}
By Lemma~\ref{lem:L-projections}, if $P=P_A$, then
\begin{equation}
\label{eq:meet-density}
       m_{x,y}(P_A)=\int_A\min\{N(x),N(y)\}\,d\mu.
\end{equation}
Indeed one minimizes pointwise by taking
$Q=P_{A\cap\{N(x)\leq N(y)\}}$.

\begin{lemma}[Intrinsic characterization]
\label{lem:beta-intrinsic}
For $x,y\in Z$, the vector $z=\beta_Z(x,y)$ is the unique vector satisfying
\begin{align}
 \|Pz\|&=m_{x,y}(P)
          &&\text{for every $L$-projection $P$,}
 \label{eq:char1}\\
 \|y\|&=\|z\|+\|y-z\|.
 \label{eq:char2}
\end{align}
Consequently every surjective linear isometry $T:Z\to Z'$ between members
of $\C_1^{\R}$ satisfies
\[
                  T\beta_Z(x,y)=\beta_{Z'}(Tx,Ty).
\]
In particular, $\beta_Z$ is independent of the chosen direct-integral or
band representation of $Z$.
\end{lemma}

\begin{proof}
Formula \eqref{eq:beta-def} gives
$N(z)=\min\{N(x),N(y)\}$, so \eqref{eq:char1} follows from
\eqref{eq:meet-density}.  Also $z(\omega)$ is a scalar multiple of
$y(\omega)$ with coefficient in $[0,1]$, and hence \eqref{eq:char2} holds.

Conversely, \eqref{eq:char1} and \eqref{eq:meet-density}, applied to every
base set, imply
\[
                       N(z)=\min\{N(x),N(y)\}
                       \quad\text{a.e.}
\]
The pointwise inequality
$N(y)\leq N(z)+N(y-z)$ has equal integrals by \eqref{eq:char2}; it is
therefore an equality almost everywhere.  Equality in the triangle
inequality of a strictly convex Hilbert space implies
$z(\omega)=t(\omega)y(\omega)$ with $0\leq t(\omega)\leq1$.  The random-norm
identity forces
\[
       t(\omega)=\frac{\min\{N(x)(\omega),N(y)(\omega)\}}{N(y)(\omega)}
\]
on $\{N(y)>0\}$, proving uniqueness.

A surjective isometry conjugates the Boolean algebras of $L$-projections,
preserves their order, and preserves every norm occurring in
\eqref{eq:meet-functional}--\eqref{eq:char2}.  It therefore sends the unique
solution for $(x,y)$ to the unique solution for $(Tx,Ty)$.
\end{proof}

\begin{lemma}[Uniform continuity]
\label{lem:beta-Lipschitz}
For every $Z\in\C_1^{\R}$,
\[
 \|\beta_Z(x,y)-\beta_Z(x',y')\|
       \leq \|x-x'\|+\|y-y'\|.
\]
\end{lemma}

\begin{proof}
For a Hilbert space, let $Q_r$ be the metric projection onto the closed ball
of radius $r$.  Metric projections onto closed convex sets are
nonexpansive, and directly
$\|Q_rv-Q_sv\|\leq|r-s|$.  Hence pointwise
\[
 \|Q_{\|u\|}v-Q_{\|u'\|}v'\|
 \leq \|u-u'\|+\|v-v'\|.
\]
Integrating gives the assertion.
\end{proof}

\begin{lemma}[Ultraproduct compatibility]
\label{lem:beta-ultra}
If $Z=\prod_{i,\U}Z_i$ is a Banach ultraproduct of members of
$\C_1^{\R}$, then
\[
 \beta_Z([x_i]_{\U},[y_i]_{\U})
       =[\beta_{Z_i}(x_i,y_i)]_{\U}.
\]
\end{lemma}

\begin{proof}
Choose $\BL_1L_2$ lattice representations and use
Proposition~\ref{prop:random-ultra}.  Put
$x=[x_i]_{\U}$, $y=[y_i]_{\U}$, and
$z=[\beta_{Z_i}(x_i,y_i)]_{\U}$.  The coordinate identities give
\begin{align*}
 N(z)&=N(x)\wedge N(y),\\
 N(y)&=N(z)+N(y-z).
\end{align*}
In the resulting $L_1(L_2)$ representation, the second equality says that
$z$ and $y-z$ have equality in the Hilbert triangle inequality in almost
every fibre.  Thus $z=t y$ for a scalar $0\leq t\leq1$, and the first
equality identifies $t$ with the radial coefficient in
\eqref{eq:beta-def}.  Therefore $z=\beta_Z(x,y)$.  Representation
independence follows from Lemma~\ref{lem:beta-intrinsic}.
\end{proof}

\section{A common-ultrapower/equalizer transfer lemma}

We isolate the model-theoretic mechanism.  No new operation is added to the
Banach-space language; the operation below is used only as an external
canonical invariant.

\begin{theorem}[Canonical-operation transfer]
\label{thm:transfer}
Let $\C$ be an isometry-invariant class of Banach spaces closed under
ultraproducts.  Suppose that every $Z\in\C$ carries a binary operation
$\beta_Z$ such that:
\begin{enumerate}[label=\textup{(\alph*)}]
\item the operations are uniformly continuous with a common modulus;
\item every surjective linear isometry between members of $\C$ preserves
      $\beta$;
\item $\beta$ commutes with Banach ultraproducts and diagonal embeddings;
\item every closed linear subspace $E\subseteq Z\in\C$ satisfying
      $\beta_Z(E,E)\subseteq E$ belongs to $\C$.
\end{enumerate}
Then $\C$ is closed under ultraroots.
\end{theorem}

\begin{proof}
Suppose that $X^{\U}$ is linearly isometric to a member of $\C$.  Transport
that structure and write
\[
                         Z=X^{\U}\in\C.
\]
Let $i=D_X:X\to Z$ be the diagonal embedding and put
\[
 e_0=D_Z:Z\longrightarrow Z^{\U},
 \qquad
 e_1=S:=i^{\U}:Z=X^{\U}\longrightarrow Z^{\U}.
\]
Both $e_0$ and $e_1$ are elementary embeddings in the pure Banach-space
language.

Expand that language by a constant symbol $c_z$ for every $z\in Z$, and
consider the two expanded structures
\[
 \mathfrak A=(Z^{\U},(e_0z)_{z\in Z}),
 \qquad
 \mathfrak B=(Z^{\U},(e_1z)_{z\in Z}).
\]
They are elementarily equivalent: the value of any formula involving
finitely many named elements is pulled back by either elementary embedding
to the same value in $Z$.  By the metric Keisler--Shelah theorem
\cite[Theorem~5.7]{BBHU2008}, there are an ultrafilter $\mathcal V$ and an
isomorphism
\[
       T:\mathfrak A^{\mathcal V}
          \longrightarrow \mathfrak B^{\mathcal V}.
\]
Write
\[
 M=(Z^{\U})^{\mathcal V},
 \qquad
 j=D_{Z^{\U}}^{\mathcal V}:Z^{\U}\longrightarrow M.
\]
The underlying space $M$ belongs to $\C$, and $T:M\to M$ is a surjective
linear isometry.  It therefore preserves $\beta$.  Moreover, because the
constants are part of the expanded language,
\[
                         Tje_0z=je_1z
                         \qquad(z\in Z).
\]
Using this identity, preservation of $\beta$ by $T$, and compatibility with
diagonal embeddings, we obtain for $a,b\in Z$
\begin{align*}
 je_1\beta_Z(a,b)
   &=Tje_0\beta_Z(a,b)\\
   &=T\beta_M(je_0a,je_0b)\\
   &=\beta_M(Tje_0a,Tje_0b)\\
   &=\beta_M(je_1a,je_1b)\\
   &=j\beta_{Z^{\U}}(e_1a,e_1b).
\end{align*}
Since $j$ is injective,
\begin{equation}
\label{eq:Sbeta}
       S\beta_Z(a,b)=\beta_{Z^{\U}}(Sa,Sb)
       \qquad(a,b\in Z).
\end{equation}

The diagonal copy of $X$ in $Z$ is the equalizer
\begin{equation}
\label{eq:equalizer}
       i(X)=\{z\in Z:D_Zz=Sz\}.
\end{equation}
Indeed, for $z=[x_k]_{\U}$ the equality $D_Zz=Sz$ is exactly
\[
       \lim_{k,\U}\|z-i(x_k)\|=0.
\]
Thus the distance from $z$ to the closed range $i(X)$ is zero, so
$z\in i(X)$; the converse is immediate.

Equations \eqref{eq:Sbeta} and \eqref{eq:equalizer}, together with the fact
that $D_Z$ preserves $\beta$, show that $i(X)$ is $\beta_Z$-closed.  By (d),
$i(X)\in\C$, and hence $X\in\C$.
\end{proof}

\begin{remark}
The common-ultrapower step is the only model-theoretic ingredient beyond
{\L}o\'s theorem.  It avoids any need to choose a specially saturated or
homogeneous ultrapower inside the class $\C$.
\end{remark}

\section{Radially closed subspaces}

This section proves condition (d) of Theorem~\ref{thm:transfer} for
$\C_1^{\R}$.

\begin{theorem}[Radial-closure representation]
\label{thm:radial-closed}
Let $Z\in\C_1^{\R}$ and let $E\subseteq Z$ be a closed real linear subspace.
If
\[
                       \beta_Z(E,E)\subseteq E,
\]
then $E\in\C_1^{\R}$.
\end{theorem}

\begin{proof}
By Proposition~\ref{prop:C1=BL}, represent $Z$ as a band in
$L_1(\Omega,\Sigma,\mu;H)$ for a fixed real Hilbert space $H$.  The radial
operation agrees with the ambient one.

\medskip
\noindent\textbf{Step 1: support decomposition.}\par
Call $x,y$ \emph{base-disjoint} when $N(x)N(y)=0$ almost everywhere.  Choose
a maximal pairwise base-disjoint family $(e_\alpha)_{\alpha\in A}$ of
nonzero elements of $E$, and put
$S_\alpha=\supp N(e_\alpha)$.  If $x\in E$, then
\[
       \beta_Z(me_\alpha,x)\longrightarrow \1_{S_\alpha}x
       \qquad(m\to\infty)
\]
in $L_1$, by dominated convergence.  Hence every restriction
$\1_{S_\alpha}x$ belongs to $E$.

For a fixed $x$, the set
\[
 A_x=\{\alpha:\|\1_{S_\alpha}x\|>0\}
\]
is countable, and the corresponding disjoint sum converges in norm.  If
$S=\bigvee_\alpha S_\alpha$ in the complete measure algebra, then for
$\alpha\notin A_x$ one has
$S_\alpha\wedge\supp N(x)=0$.  Since finite meets distribute over arbitrary
joins in a complete Boolean algebra,
$N(x)=0$ on $S\setminus\bigvee_{\alpha\in A_x}S_\alpha$.  Hence
\[
              \1_Sx=\sum_{\alpha\in A_x}\1_{S_\alpha}x\in E.
\]
The residual part $\1_{S^c}x$ belongs to $E$ and, if nonzero, would be
base-disjoint from every $e_\alpha$, contradicting maximality.  Therefore
\begin{equation}
\label{eq:E-decomp}
       E=\left(\bigoplus_{\alpha\in A}E_\alpha\right)_1,
       \qquad E_\alpha=\1_{S_\alpha}E.
\end{equation}
It is enough to treat one summand.  We may thus assume that $E$ contains an
element $e$ with
\[
       h=N(e)>0\quad\text{almost everywhere on the common support of }E.
\]

\medskip
\noindent\textbf{Step 2: normalization.}
Restrict to that support and set $d\nu=h\,d\mu$.  Notice that
$\nu(\Omega)=\|e\|<\infty$.  Define
\[
        T:E\longrightarrow L_1(\Omega,\Sigma,\nu;H),
        \qquad Tx=h^{-1}x.
\]
This is a linear isometry and it intertwines radial truncation.  Put
$Y=T(E)$ and $u=T(e)$.  Then $N(u)=1$ almost everywhere and
$\beta(Y,Y)\subseteq Y$.

\medskip
\noindent\textbf{Step 3: extraction of a scalar sub-$\sigma$-algebra.}
Define
\[
       A_0=\{f\in L_1(\Omega,\Sigma,\nu):fu\in Y\}.
\]
The map $f\mapsto fu$ is an isometry, so $A_0$ is a closed real linear
subspace of scalar $L_1$ and $1\in A_0$.  For $f\in A_0$,
\[
       \beta(fu,mu)=\min\{|f|,m\}u\longrightarrow |f|u.
\]
Thus $A_0$ is closed under modulus and hence is a closed vector sublattice.
For $y\in Y$,
\[
       \beta(y,mu)=\min\{N(y),m\}u\longrightarrow N(y)u,
\]
so
\begin{equation}
\label{eq:Ny-in-A}
                       N(y)\in A_0\qquad(y\in Y).
\end{equation}

Let
\[
       \mathcal F=\{B\in\Sigma:\1_B\in A_0\},
\]
modulo null sets.  Because $A_0$ is a closed sublattice containing $1$,
$\mathcal F$ is a $\sigma$-algebra: complements and finite intersections
follow from the lattice operations, and countable unions follow by monotone
$L_1$ convergence.  If $f\in A_0$ and $t\in\mathbb Q$, then
\[
       1\wedge n(f-t)^+\longrightarrow \1_{\{f>t\}}
       \quad\text{in }L_1,
\]
so $f$ is $\mathcal F$-measurable.  Conversely, $\mathcal F$-simple
functions belong to $A_0$ and are dense in $L_1(\mathcal F)$.  Hence
\begin{equation}
\label{eq:A=L1F}
                       A_0=L_1(\Omega,\mathcal F,\nu).
\end{equation}

\medskip
\noindent\textbf{Step 4: module structure.}
If $B\in\mathcal F$, then $\1_Bu\in Y$, and for every $y\in Y$,
\[
       \beta(m\1_Bu,y)\longrightarrow \1_By
       \qquad(m\to\infty).
\]
Therefore $Y$ is closed under multiplication by indicators in $\mathcal F$,
and then, by uniform approximation, under all multipliers in
$L_\infty(\mathcal F)$.  By \eqref{eq:Ny-in-A} and
\eqref{eq:A=L1F}, every random norm $N(y)$ is $\mathcal F$-measurable.

After completing both $(\Omega,\Sigma,\nu)$ and
$(\Omega,\mathcal F,\nu)$, they are finite and therefore localizable.
Raynaud's module representation lemma
\cite[Lemma~4.4]{Raynaud2004} now applies and gives pairwise almost disjoint
sets $(\Omega_j)_{j\in J}$ in $\mathcal F$, real Hilbert spaces $H_j$, and a
random-norm-preserving linear isometry
\[
       Y\cong
       \left(\bigoplus_{j\in J}
       L_1(\Omega_j,\mathcal F|_{\Omega_j},\nu|_{\Omega_j};H_j)
       \right)_1.
\]
Thus the normalized component belongs to $\C_1^{\R}$.  Undoing $T$ and
using the $\ell_1$-decomposition \eqref{eq:E-decomp} proves
$E\in\C_1^{\R}$.
\end{proof}

\section{Proof of the main theorem}

\begin{proof}[Proof of Theorem~\ref{thm:main}]
By Proposition~\ref{prop:C1=BL}, it is enough to prove that
$\C_1^{\R}$ is axiomatizable.  This class is invariant under surjective
linear isometries and is closed under ultraproducts by
Proposition~\ref{prop:random-ultra}.  Lemmas
\ref{lem:beta-intrinsic}, \ref{lem:beta-Lipschitz}, and
\ref{lem:beta-ultra} show that its radial operation is canonical, uniformly
continuous, and compatible with ultraproducts.  Theorem
\ref{thm:radial-closed} says that every closed radial-invariant subspace of a
member is again in the class.  The canonical-operation transfer theorem
therefore implies closure under ultraroots.

A class of Banach structures is axiomatizable exactly when it is closed under
isometries, ultraproducts, and ultraroots
\cite[Proposition~5.14]{BBHU2008}.  Hence $(\BL_1L_2)^B$ is axiomatizable.
\end{proof}

\begin{corollary}
Raynaud's Question~6.17 has a positive answer at the real endpoint $q=2$.
\end{corollary}


\section{Why the stable-variable counterargument does not apply}

A proposed negative argument uses the isometric stable-variable embedding
$\ell_q^n\hookrightarrow L_1$ for $1<q\leq2$ and concludes that
$L_1$ and $L_1(\ell_q^n)$ have isometric ultrapowers.  The first step is
correct, but the conclusion is not.  Heinrich's theorem characterizes finite
representability by an isometric \emph{embedding} into an ultrapower
\cite{Heinrich1980}; mutual finite representability yields embeddings in both
directions, not a surjective isometry between ultrapowers.  Hence that
argument does not contradict Theorem~\ref{thm:main}.

\section*{Verification notes}

The proof has three external inputs:
\begin{enumerate}[leftmargin=2.2em]
\item the mixed-norm random-norm representation of ultraproducts
      (Proposition~\ref{prop:random-ultra});
\item the metric Keisler--Shelah theorem for the two constant expansions
      in the proof of Theorem~\ref{thm:transfer};
\item Raynaud's module representation lemma
      \cite[Lemma~4.4]{Raynaud2004}, used with $p=1$.
\end{enumerate}
All other ingredients, including the intrinsic characterization of $\beta$,
the equalizer argument, and the extraction of $\mathcal F$, are proved in
the note.  Since the endpoint conclusion appears to be new, these three interfaces---and
in particular the passage from the intrinsic radial operation to the
common-ultrapower equalizer---are the natural points for independent expert
verification.

\begin{thebibliography}{99}

\bibitem{BBHU2008}
I.~Ben Yaacov, A.~Berenstein, C.~W.~Henson, and A.~Usvyatsov,
\emph{Model theory for metric structures},
in Model Theory with Applications to Algebra and Analysis, Vol.~2,
London Math. Soc. Lecture Note Ser. 350, Cambridge Univ. Press, 2008,
315--427.

\bibitem{Heinrich1980}
S.~Heinrich,
\emph{Ultraproducts in Banach space theory},
J. Reine Angew. Math. \textbf{313} (1980), 72--104.


\bibitem{HensonRaynaud2007}
C.~W.~Henson and Y.~Raynaud,
\emph{On the theory of $L_p(L_q)$-Banach lattices},
Positivity \textbf{11} (2007), 201--230.


\bibitem{HensonRaynaud2011}
C.~W.~Henson and Y.~Raynaud,
\emph{Quantifier elimination in the theory of $L_p(L_q)$-Banach lattices},
J. Log. Anal. \textbf{3} (2011), Paper 11, 1--29.

\bibitem{LevyRaynaud1984}
M.~L\'evy and Y.~Raynaud,
\emph{Ultrapuissances des espaces $L^p(L^q)$},
C. R. Acad. Sci. Paris S\'er. I Math. \textbf{299} (1984), 81--84;
see also S\'eminaire d'analyse fonctionnelle 1983/84, 69--79.

\bibitem{Raynaud2004}
Y.~Raynaud,
\emph{The range of a contractive projection in $L_p(H)$},
Rev. Mat. Complut. \textbf{17} (2004), no.~2, 485--512.

\bibitem{Raynaud2018}
Y.~Raynaud,
\emph{New axiomatizable classes of Banach spaces via
disjointness-preserving isometries},
Positivity \textbf{22} (2018), 301--339.

\end{thebibliography}

\end{document}
